\section{Introduction}\label{sec:intro}
Planners and schedulers who use MIP solvers rarely wait for a proof of optimality; they need a good feasible solution
within a fixed time. How soon a solver finds one is measured by the primal integral $P(T)$~\citep{berthold2013primal},
the gap to the best known value averaged over the time budget $T$, which rewards finding good solutions early.

GPUs now sit beside CPU solvers on small machines, and they suit population search: many candidate solutions can be
evaluated and repaired at once. NVIDIA cuOpt already runs a GPU population of solutions with recombination and local
search beside its own branch-and-bound~\citep{corduk2025gpuprimal,corduk2025evoframework,nvidia2026cuoptsrc}. In an
earlier study on the same machine, simply running cuOpt beside HiGHS or SCIP lowered the primal integral by 21--30\,\%
at 60\,s, while a CPU evolutionary algorithm exchanging solutions with SCIP gave no measurable gain and lost to a second
copy of SCIP~\citep{orbay2026optopt}.

I-TREA II tests the changes that study named: a population that lives on the GPU, longer budgets (300\,s and
1{,}800\,s), and instances on which the solver's own heuristics struggle. The evolutionary algorithm (EA) runs beside
an exact solver --- SCIP as the primary partner, HiGHS as a replication partner --- and the two exchange solutions in
both directions: the solver's incumbent, very good solutions from its pool, its proven bound and its root LP solution
enter the population, and the population's better incumbents enter the solver as new starting points. From the first
I-TREA~\citep{orbay2011itrea} it keeps one principle: good infeasible solutions --- infeasible solutions whose objective
is better than the best feasible one --- are kept in the population, here as a capped share (Rank~0), beside a
protected share of the best feasible solutions (Rank~1).

\paragraph{Question.} The claim under test is that a GPU-based evolutionary metaheuristic coupled with an exact solver
can beat either one alone. On hard MIP instances, does the coupled system find good feasible solutions sooner than
\emph{the solver alone} and \emph{the EA alone}? Both comparisons must succeed for the claim to hold. Running two methods
side by side can beat either one simply by being two, so we also compare the coupling with the same two methods running
without exchange, and with the solver racing cuOpt on the same GPU (Section~\ref{sec:design}).

\paragraph{Closest prior work.} None of the parts is new. Incumbents were exchanged in both directions between a CPU
heuristic and a branch-and-bound solver running in parallel two decades ago~\citep{nwana2005cooperative,
puchinger2005cooperating}. ReXi runs local-search and large-neighbourhood workers beside SCIP workers over a shared
solution pool~\citep{mexi2026rexi}; CHAP pools the feasible and infeasible candidates of GPU and CPU
heuristics~\citep{tjusila2026chap}; and inside cuOpt a GPU population with recombination, which keeps infeasible
members and protects the best feasible one, runs beside cuOpt's own CPU branch-and-bound with solutions flowing both
ways~\citep{corduk2025evoframework,nvidia2026cuoptsrc}. Reserving a share of the population for infeasible solutions
goes back at least to IDEA~\citep{singh2011thesis}. What this study puts together is a GPU population coupled with an
\emph{external, general} MIP solver through two-way exchange, with a capped share of infeasible solutions whose
objective beats the incumbent, and with the solver's proven bound removing those that can never become feasible. We did
not find this combination in our search, which is a negative search result, not a priority claim
(Section~\ref{sec:related}).

\paragraph{Status of this draft.} The hypotheses, arms, instance-pool rules, metrics and analysis were fixed before any
experiment and are summarised in Section~\ref{sec:design}. No experiment has been run; Section~\ref{sec:results} is a
placeholder. \pending{Contributions to be stated after the results; no claim is made here.}
